Aux sports d'hiver, un enfant sur une luge part, sans vitesse, du sommet d'une
pente de longueur $L=\qty{30.0}{\m}$ inclinée de $\alpha=\ang{30}$ sur
l'horizontale.

Au bas de cette pente, il aborde une autre pente inclinée de l'angle
$\beta=\ang{15}$ sur l'horizontale.

\begin{minipage}{\linewidth}
    \begin{figure}[H]
        \centering
        \includegraphics[width=\textwidth]{dbbbd740-7087-40eb-bf09-fa10c8492051-3}
    \end{figure}
\end{minipage}

Quelle distance $L^{\prime}$ parcourt-il le long de cette deuxième pente avant
de s'arrêter~?

On négligera tous les frottements. On admettra en outre que la cassure de la
pente au point le plus bas de la trajectoire ne modifie pas la valeur de la
vitesse.
